% !TeX TS-program = xelatex

\documentclass[compress,aspectratio=169]{beamer}

\usepackage{presentationtemplate}
\usepackage[fontsize=\footnotesize, askip=3mm]{terminal}
\usepackage[fontsize=\footnotesize, linenosfontsize=\tiny, askip=3mm]{mylisting}
\usepackage{myitemize}

\title{Время хранения и время жизни объектов}
\subtitle{Storage duration, время жизни и связывание в C}

\begin{document}

\frame[plain]{\titlepage}

\begin{frame}

    \frametitle{Два независимых вопроса о переменной}

    У объявления переменной две характеристики,
    которые часто путают, — и они не определяют
    друг друга.

    \begin{myitemize}

        \item scope (область видимости) — где в тексте
            программы имя ссылается на именно этот
            объект; свойство времени компиляции.

        \item storage duration (время хранения) — как
            долго объект занимает память; определяет
            lifetime (время жизни) во время выполнения.

        \item Пример разрыва: \texttt{static}-переменная
            внутри функции — имя видно только в функции
            (scope), а объект живёт всю программу
            (duration).

    \end{myitemize}

\end{frame}

\begin{frame}

    \frametitle{Lifetime и storage duration}

    \begin{myitemize}

        \item lifetime (время жизни) — часть времени
            выполнения, в течение которой объект занимает
            память: от создания до уничтожения.

        \item storage duration — свойство времени
            компиляции, которое задаёт минимальное время
            жизни объекта.

        \item Виды в C: \texttt{automatic}, \texttt{static},
            \texttt{dynamic}, \texttt{thread}.

        \item Thread storage duration (C11,
            \texttt{\_Thread\_local}) — упоминаем одним
            словом; отдельно не рассматриваем.

    \end{myitemize}

\end{frame}

\begin{frame}[fragile]

    \frametitle{Automatic: время жизни объекта}

    \begin{myitemize}

        \item Automatic-объект создаётся при входе в блок
            и уничтожается при выходе из него.

        \item При повторном входе в блок создаётся новый
            объект — предыдущий уже не существует.

    \end{myitemize}

    \begin{myinplacelisting}[minted language=c]
int main(void)
{
    char c0 = 0; /* c0: begins */
    (void)c0;
    {
        char c1 = 0; /* c1: begins */
        (void)c1;
    } /* c1: ends */
    return 0;
} /* c0: ends */
    \end{myinplacelisting}

\end{frame}

\begin{frame}

    \frametitle{Automatic: LIFO-стек}

    \begin{columns}[T]

        \begin{column}{0.54\textwidth}

            \begin{myitemize}

                \item На практике automatic storage duration
                    реализуется LIFO-стекком (stack-based
                    memory allocation).

                \item Вызов функции кладёт на стек кадр
                    (stack frame) со своими
                    automatic-объектами; возвращение из
                    функции снимает кадр.

                \item Отсюда ограничение: глубина рекурсии
                    ограничена объёмом стека
                    (stack overflow).

            \end{myitemize}

        \end{column}

        \begin{column}{0.44\textwidth}

            \begin{center}

                \begin{tikzpicture}[
                    bytenode/.style={
                        rectangle,
                        thick,
                        minimum width=9mm,
                        minimum height=6mm,
                    },
                ]
                    \node[bytenode, fill=green!10,
                        label={[label distance=1.5mm]left:{\footnotesize\texttt{0x7ffccbf93240}}},
                        label={[label distance=1.5mm]right:\footnotesize\textit{свободна}}]
                        (b0) at (0, 0) {};
                    \node[bytenode, fill=yellow!10, above=0 of b0,
                        label={[label distance=1.5mm]left:{\footnotesize\texttt{0x7ffccbf9323f}}}]
                        (b1) {\texttt{c0}};
                    \node[bytenode, fill=red!10, above=0 of b1,
                        label={[label distance=1.5mm]left:{\footnotesize\texttt{0x7ffccbf9323e}}}]
                        (b2) {\texttt{c1}};
                \end{tikzpicture}

                \small
                \texttt{c1} создан над \texttt{c0} (по меньшему
                адресу): LIFO — при выходе из вложенного блока
                \texttt{c1} уничтожается раньше.

            \end{center}

        \end{column}

    \end{columns}

\end{frame}

\begin{frame}[fragile]

    \frametitle{Automatic в C: неинициализированная переменная}

    \begin{myitemize}

        \item Indeterminate value — это не ноль: в памяти
            может быть что угодно.

        \item Инициализируйте automatic-переменные
            при объявлении.

    \end{myitemize}

    \begin{columns}[T]

        \begin{column}{0.44\textwidth}

            \myinputlisting[minted language=c]
                {presentations/03-storage-duration/}
                {uninit.c}

        \end{column}

        \begin{column}{0.54\textwidth}

            \begin{terminalwindow}[minted options={fontsize=\scriptsize, breaklines, breakanywhere, escapeinside=!!, breaksymbolleft=}]
!!\shellcommand{gcc -Wall uninit.c}!!
uninit.c: In function ‘main’:
uninit.c:6:5: warning: ‘i’ is used uninitialized [-Wuninitialized]
    6 |     printf("%d\n", i);
      |     ^~~~~~~~~~~~~~~~~
uninit.c:5:9: note: ‘i’ was declared here
    5 |     int i;
      |         ^
            \end{terminalwindow}

        \end{column}

    \end{columns}

\end{frame}

\begin{frame}[fragile]

    \frametitle{Dynamic: malloc и free}

    \begin{columns}[T]

        \begin{column}{0.44\textwidth}

            \begin{myitemize}
                \footnotesize

                \item Куча (heap) — регион памяти,
                    отдельный от стека, без
                    LIFO-порядка стека.

                \item Время жизни объекта: от успешного
                    \texttt{malloc} до \texttt{free}
                    (или до завершения программы).

                \item В C — только \texttt{malloc}/
                    \texttt{free} (нет \texttt{new}/
                    \texttt{delete}); \texttt{free(p)} —
                    конец времени жизни объекта,
                    указатель \texttt{p} не уничтожается.

            \end{myitemize}

        \end{column}

        \begin{column}{0.54\textwidth}

            \myinputlisting[minted language=c, minted options={fontsize=\tiny, breaklines, breakanywhere, linenos, escapeinside=||, numbersep=2pt}]
                {presentations/03-storage-duration/}
                {malloc-heap.c}

        \end{column}

    \end{columns}

\end{frame}

\begin{frame}

    \frametitle{Dynamic: правила работы с кучей}

    \begin{myitemize}

        \item \texttt{malloc} возвращает \texttt{void*} —
            в C \texttt{void*} неявно приводится к любому
            указателю на объект, каст не нужен.

        \item При нехватке памяти \texttt{malloc}
            возвращает \texttt{NULL} — результат
            обязательно проверяйте до использования.

        \item Каждому \texttt{malloc} — ровно один
            \texttt{free}; использование памяти после
            \texttt{free} и повторный \texttt{free} —
            неопределённое поведение.

        \item Размер выделяется в байтах:
            \texttt{n * sizeof *p}.

    \end{myitemize}

\end{frame}

\begin{frame}[fragile]

    \frametitle{Static: время хранения и нулевая инициализация}

    \begin{columns}[T]

        \begin{column}{0.44\textwidth}

            \myinputlisting[minted language=c]
                {presentations/03-storage-duration/}
                {static-objects.c}

        \end{column}

        \begin{column}{0.54\textwidth}

            \begin{myitemize}
                \footnotesize

                \item \texttt{int g;} (global): живёт всю
                    программу и инициализируется нулём
                    автоматически: скаляры — 0, указатели
                    — \texttt{NULL}.

                \item \texttt{static int h;} (file scope):
                    то же время хранения и нулевая
                    инициализация; плюс внутреннее
                    связывание (слайд 11).

                \item \texttt{static int c;} в функции: имя
                    видно только в этой функции, объект
                    живёт всю программу и начинается с 0.

                \item Для сравнения: automatic-переменная
                    без инициализатора — неопределённое
                    значение (слайд 6).

            \end{myitemize}

        \end{column}

    \end{columns}

\end{frame}

\begin{frame}[fragile]

    \frametitle{Static: счётчик между вызовами}

    \begin{columns}[T]

        \begin{column}{0.6\textwidth}

            \myinputlisting[minted language=c, minted options={fontsize=\scriptsize, breaklines, breakanywhere, linenos, escapeinside=||, numbersep=2pt}]
                {presentations/03-storage-duration/}
                {static-counter.c}

        \end{column}

        \begin{column}{0.38\textwidth}

            \begin{terminalwindow}
!!\shellcommand{gcc static-counter.c}!!
!!\shellcommand{./a.out}!!
1
2
3
            \end{terminalwindow}

        \end{column}

    \end{columns}

\end{frame}

\begin{frame}[fragile]

    \frametitle{Linkage: internal и external}

    Связывание (linkage) определяет, можно ли обращаться
    к имени объекта (или функции) из другой единицы
    трансляции.

    \begin{myitemize}
        \footnotesize

        \item Единица трансляции — один .c-файл
            после препроцессирования.

        \item External linkage — по умолчанию у
            глобальных объектов и функций: имя доступно
            из других файлов (через \texttt{extern}).

        \item Internal linkage — глобальный объект
            с \texttt{static} в начале файла: имя
            существует только в этом файле (защита от
            конфликтов имён).

        \item В C нет namespaces: конфликты имён
            между файлами решаются префиксами-конвенциями
            (см. 02-c-vs-cpp).

    \end{myitemize}

    \begin{myinplacelisting}[minted language=c]
int x;         /* external linkage */
static int y;  /* internal linkage */
    \end{myinplacelisting}

\end{frame}

\begin{frame}[fragile]

    \frametitle{Extern: одна переменная в нескольких файлах}

    \small
    Объект определяется один раз (\texttt{foo.c});
    \texttt{extern} в заголовке — объявление без
    выделения памяти, поэтому все файлы видят один
    и тот же объект.

    \normalsize

    \begin{columns}[T]

        \begin{column}{0.48\textwidth}

            \myinputlisting[minted language=c]
                {presentations/03-storage-duration/extern/}
                {foo.h}

            \myinputlisting[minted language=c]
                {presentations/03-storage-duration/extern/}
                {foo.c}

        \end{column}

        \begin{column}{0.48\textwidth}

            \myinputlisting[minted language=c, minted options={fontsize=\scriptsize, breaklines, breakanywhere, linenos, escapeinside=||, numbersep=2pt}, after skip=0mm]
                {presentations/03-storage-duration/extern/}
                {main.c}

            \begin{terminalwindow}[minted options={fontsize=\scriptsize, breaklines, breakanywhere, escapeinside=!!, breaksymbolleft=}]
!!\shellcommand{gcc main.c foo.c}!!
!!\shellcommand{./a.out}!!
1
            \end{terminalwindow}

        \end{column}

    \end{columns}

\end{frame}

\begin{frame}

    \frametitle{Итоговая таблица}

    \begin{center}

        \scriptsize
        \setlength{\tabcolsep}{4pt}
        \begin{tabular}{lll}
            \hline
            \textbf{storage duration} & \textbf{пример} & \textbf{время жизни} \\
            \hline
            automatic & \texttt{int x;} в функции &
                вход в блок $\rightarrow$ выход из блока \\
            static & \texttt{int g;}, \texttt{static int c;} &
                запуск $\rightarrow$ завершение программы \\
            dynamic & \texttt{p = malloc(n)} &
                \texttt{malloc} $\rightarrow$ \texttt{free} \\
            thread & \texttt{\_Thread\_local int x;} (C11) &
                старт $\rightarrow$ конец потока \\
            \hline
        \end{tabular}

    \end{center}

    \vspace{1em}

    \begin{myitemize}

        \item Нулевая инициализация по умолчанию: static
            (глобальные и \texttt{static}) — да; automatic
            и dynamic — нет (неопределённое значение).

    \end{myitemize}

\end{frame}

\begin{frame}

    \frametitle{Выводы}

    \begin{myitemize}

        \item Scope и storage duration — независимые
            свойства: где видно имя — не значит, сколько
            живёт объект.

        \item Automatic без инициализатора в C —
            indeterminate value: чтение до записи — UB
            (предупреждение \texttt{-Wuninitialized}).

        \item Dynamic: \texttt{malloc}/\texttt{free},
            проверка \texttt{NULL}, один \texttt{malloc}
            — один \texttt{free}.

        \item Static: живёт всю программу, инициализируется
            нулём; \texttt{static} в file scope — internal
            linkage.

        \item Глобальные объекты с external linkage
            разделяются между файлами через
            \texttt{extern}.

    \end{myitemize}

\end{frame}

\end{document}
